\paragraph{A limit of the valuation-one argument.}
The valuation hypothesis in CX-N14 cannot simply be deleted from
its particular noncyclicity proof.
For an explicit example, take \(p=17\), \(r=3\), and
\(A=\mathbb F_{17}^3\).
The polynomial \(T^2+T+1\) has no root over \(\mathbb F_{17}\),
since a root would have order three in its multiplicative group
of order sixteen. Thus
\[
 D=\mathbb F_{17}[T]/(T^2+T+1)\cong\mathbb F_{17^2},
 \qquad \zeta=T\bmod(T^2+T+1).
\]
Here \(s=\operatorname{ord}_3(17)=2\) and
\(v_3(17^2-1)=v_3(288)=2\).
The cyclic group \(D^\times\) contains an element \(\eta\) of
order nine with \(\eta^3=\zeta\): start with a generator of its
order-nine subgroup and replace it by its inverse if necessary.
Set
\[
 x=(\eta,\zeta\eta,\zeta^2\eta)\in
 A\otimes_{\mathbb F_{17}}D=D^3,
 \qquad
 \tau(a_0,a_1,a_2)=(a_1,a_2,a_0).
\]
Then \(\tau(x)=\zeta x\), \(x\) has order nine, and
\(x^3=(\zeta,\zeta,\zeta)\).
The nonidentity cyclotomic automorphism \(\rho_2\) of \(D\)
is its seventeenth-power Frobenius.
On every element of order dividing nine, this acts by exponent
\(17\equiv-1\pmod9\), which is also
\(\omega_2(2)=2^3\equiv-1\pmod9\).
It follows coordinate by coordinate that
\(\rho_2(x)=x^{\omega_2(2)}\); hence \(x\) is in the required
Teichm\"uller subgroup.
Nevertheless \(\langle x,\zeta\rangle=\langle x\rangle\)
is cyclic of order nine.
This example satisfies the eigenunit properties used from
Proposition~4.6 but does not force that pair to be noncyclic.
It makes no assertion about which valid eigenunit a particular
implementation returns, or about the impossibility of a different
Dirichlet-only algorithm.
\hfill\(\square\)
